%%%%%%%%%%%%%%%%%%%%
% INTRODUCCION
%%%%%%%%%%%%%%%%%%%%

La arquitectura de un sistema de software no es una propiedad estática: se moldea, se negocia y, con frecuencia, se deteriora a medida que el código evoluciona. Cuando las decisiones de implementación se apartan de la estructura prevista sin que nadie lo registre, aparece la deriva arquitectónica, un distanciamiento paulatino entre la arquitectura deseada y la que efectivamente quedó plasmada en el código fuente \citep{PerryWolf1992, Li2022}. Se trata de un fenómeno silencioso, acumulativo y costoso, que rara vez se manifiesta de golpe; por el contrario, avanza mediante pequeñas concesiones —una dependencia que cruza una capa, un atajo entre módulos, una clase que concentra demasiadas responsabilidades— que, tomadas de forma aislada, parecen inofensivas.

Las consecuencias de ese deterioro han sido ampliamente documentadas. La erosión arquitectónica incrementa de forma sostenida el esfuerzo de mantenimiento, encarece la incorporación de nuevas funcionalidades y reduce la capacidad del equipo para predecir el impacto de sus cambios \citep{Pigazzini2021}. Estudios empíricos recientes confirman que ciertos \textit{smells} arquitectónicos correlacionan de manera negativa con la mantenibilidad del software \citep{Jolak2025}, y que los propios desarrolladores perciben la deriva como una fuente persistente de fricción en su trabajo cotidiano \citep{Anthony2024}. El problema no es menor: en un estudio con desarrolladores de la industria, replicado sobre un segundo conjunto de datos independiente, \citet{Besker2019} estimaron que la deuda técnica consume en promedio el $23\%$ del tiempo de desarrollo, y la revisión sistemática de \citet{Li2022} identifica entre las consecuencias de la erosión el deterioro de la calidad y una evolución cada vez más rígida y frágil.

El escenario se agudiza en las pequeñas y medianas empresas (PyMES) de software. A diferencia de las grandes organizaciones, estas compañías suelen operar con equipos reducidos, ciclos de entrega ajustados y una documentación arquitectónica escasa, desactualizada o, en muchos casos, inexistente \citep{Hamed2023}. Carecen, además, de los recursos humanos y económicos para desplegar las plataformas de verificación sofisticadas que sí están al alcance de corporaciones con mayor madurez. El resultado es una paradoja preocupante: precisamente donde la gobernanza técnica sería más rentable, es donde menos se practica. Frente a este panorama, resulta pertinente preguntarse si es posible trasladar los beneficios de la verificación arquitectónica a un entorno de recursos limitados, mediante un mecanismo ligero, automatizado y que no dependa de documentación externa.

Este trabajo propone responder esa inquietud a través de un mecanismo que combina el análisis estático de grafos de dependencia con la ejecución de \textit{Fitness Functions}. Estas últimas, popularizadas en la práctica de la ingeniería de software, consisten en verificaciones objetivas y automatizables que expresan, en forma de criterios cuantificables, las propiedades que la arquitectura debe preservar \citep{PaulWang2019, FordParsonsKua2017, Ford2022}. En lugar de confiar en un documento de arquitectura que probablemente esté desactualizado, el mecanismo infiere la estructura real del sistema a partir de su propio código, la contrasta contra un conjunto acotado de reglas estructurales y emite un diagnóstico sobre el estado de conformidad del proyecto. De esta forma, la gobernanza arquitectónica deja de ser un ejercicio documental y se convierte en una práctica continua, verificable y de bajo costo.

\section{Antecedentes y contexto}
La preocupación por mantener la coherencia entre el diseño previsto y el código construido no es reciente. Desde que la ingeniería de software reconoció a la arquitectura como una disciplina con identidad propia, a comienzos de la década de 1990, quedó establecido que la estructura de un sistema condiciona de manera decisiva sus atributos de calidad: rendimiento, mantenibilidad, seguridad y capacidad de evolución \citep{PerryWolf1992, Bass2012}. La arquitectura dejó de entenderse como un dibujo accesorio y pasó a concebirse como el andamiaje sobre el que se sostienen todas las decisiones posteriores; de ahí que su deterioro no sea un asunto cosmético, sino una amenaza concreta contra la viabilidad del producto.

Paralelamente, la evidencia empírica consolidó una constatación incómoda: los sistemas reales cambian, y cambian mucho. Las leyes de evolución del software enunciadas por \citet{Lehman1980} anticiparon que todo programa que resulte útil a su comunidad se verá sometido a una presión continua de modificación, y que, a medida que crece, su estructura tiende a degradarse a menos que se invierta un esfuerzo explícito en preservarla. La deriva arquitectónica es, en buena medida, la manifestación contemporánea de esa ley: el costo de no intervenir se paga, tarde o temprano, en forma de complejidad estructural y deuda acumulada.

El concepto de deuda técnica, acuñado por \citet{Cunningham1992} para describir las consecuencias de aplazar decisiones de diseño, proporcionó el vocabulario con el que la industria pudo nombrar el fenómeno. Desde entonces, la gestión de esa deuda se convirtió en un campo de estudio activo, y su variante estructural —la deuda técnica arquitectónica— en uno de sus frentes más difíciles de atender, precisamente porque sus síntomas no se manifiestan en una línea de código aislada, sino en la relación entre componentes que, por separado, parecen saludables.

En el ámbito de las pequeñas y medianas empresas, esta problemática se cruza con otra restricción. La literatura sobre transformación digital de PyMES coincide en señalar que estas organizaciones necesitan soluciones de bajo costo, de adopción incremental y que no exijan especialistas dedicados \citep{Patricio2025, Jamil2026}. Los marcos de gobernanza pensados para corporaciones resultan, por su peso y su demanda de roles, poco realistas para equipos pequeños. Esta tensión entre la necesidad de control estructural y la escasez de recursos constituye el terreno fértil sobre el que se levanta la presente propuesta.

\section{Definición del problema}

\subsection{Planteamiento del problema}
El punto de partida es una situación concreta y, a la vez, extendida: los sistemas de software de las PyMES crecen sin que exista un mecanismo que vigile, de manera continua, la coherencia entre lo que se pretendía construir y lo que realmente se construyó. Los síntomas son reconocibles para cualquier profesional del área: dependencias que cruzan fronteras que no debían cruzarse, módulos que se acoplan de forma imprevista, componentes que se vuelven imprescindibles para todo el sistema y una dificultad creciente para estimar cuánto costará el siguiente cambio. Las causas son múltiples, pero se concentran en la ausencia de documentación arquitectónica viva, la presión por entregar con rapidez y la falta de herramientas de verificación asequibles para equipos pequeños \citep{Hamed2023, Klein2022}.

Si la situación se mantiene, los efectos negativos se profundizan. La deuda técnica se acumula en forma de \textit{smells} arquitectónicos que, según la evidencia disponible, terminan comprometiendo la mantenibilidad y la longevidad de los productos \citep{SasAvgeriou, Jolak2025}. Cada desviación no detectada encarece las siguientes: el equipo invierte tiempo en comprender un código cuya estructura real ya no se corresponde con ninguna representación mental compartida. La organización, en consecuencia, ve erosionada su capacidad de respuesta ante el mercado, justo cuando más la necesita.

\subsection{Sistematización del problema}
A partir de lo anterior, la investigación se orienta mediante una pregunta central, formulada de manera abierta para que actúe como guía y no como una búsqueda de una única causa:

\textit{¿Es posible diseñar y evaluar un mecanismo automatizado, basado en análisis estático, grafos de dependencia y \textit{Fitness Functions}, que permita detectar tempranamente derivas arquitectónicas y clasificar el estado de conformidad estructural de un sistema de software para PyMES sin depender de documentación arquitectónica externa?}

De esta pregunta se desprenden interrogantes subsidiarios que delimitan las etapas del trabajo: qué propiedades estructurales deben vigilarse para que el diagnóstico sea útil y acotado; de qué manera puede obtenerse, a partir del código fuente, un grafo de dependencias lo bastante fiel para soportar la evaluación; cómo debe traducirse el resultado de las verificaciones en una clasificación del estado arquitectónico que sea comprensible para un equipo de desarrollo; y, finalmente, cómo evidenciar, mediante un experimento controlado, que el mecanismo efectivamente detecta las derivas que se le inducen.

\subsection{Hipótesis de trabajo}
\label{sec:hipotesis-intro}
El anteproyecto formuló la hipótesis que orienta la evaluación: si se implementa un mecanismo automatizado que extraiga dependencias del código, ejecute \textit{Fitness Functions} sobre un grafo estructural y clasifique los resultados mediante umbrales explícitos de conformidad, entonces será posible (i) detectar tempranamente derivas arquitectónicas relevantes en un proyecto de software de PyME, (ii) diferenciar de manera objetiva un estado aceptable de uno problemático y (iii) producir evidencia suficiente para apoyar decisiones de mantenimiento y refactorización, sin depender de documentación arquitectónica externa. Los tres componentes se contrastan por separado, con la evidencia que corresponde a cada uno y declarando lo que no quedó demostrado, en la sección~\ref{sec:hipotesis}.

\section{Objetivos del proyecto}

\subsection{Objetivo general}
Desarrollar y evaluar experimentalmente un mecanismo automatizado para la detección de derivas arquitectónicas, basado en grafos de dependencia y en la ejecución de \textit{Fitness Functions}, que permita clasificar el estado de conformidad estructural de proyectos de software de PyMES sin depender de documentación arquitectónica externa.

\subsection{Objetivos específicos}
\begin{itemize}
    \item \textbf{OBJ1.} Formular un catálogo acotado de \textit{Fitness Functions} que traduzcan, en criterios verificables, las propiedades estructurales esenciales que deben preservarse en un sistema de software (integridad de capas, ausencia de ciclos y control de la concentración de dependencias salientes).
    \item \textbf{OBJ2.} Implementar un componente de análisis estático capaz de construir, a partir del código fuente, un grafo de dependencias entre paquetes que sirva de insumo común para todas las verificaciones.
    \item \textbf{OBJ3.} Construir un motor de validación que ejecute las \textit{Fitness Functions} sobre el grafo, agregue sus resultados en una Tasa de Conformidad (\textit{CR}) y clasifique el estado arquitectónico del sistema en niveles interpretables.
    \item \textbf{OBJ4.} Ejecutar un experimento controlado sobre un caso de estudio para comprobar que el mecanismo detecta las derivas inducidas y clasifica correctamente tanto una implementación de referencia como una versión alterada.
\end{itemize}

\section{Delimitaciones y alcances}
El trabajo se circunscribió a un conjunto de fronteras definidas de antemano, con el propósito de hacer viable la propuesta en el tiempo y con los recursos disponibles. En primer lugar, el mecanismo se orientó exclusivamente al ecosistema Java, aprovechando la madurez de las herramientas de análisis estático disponibles para esa plataforma; no se abordaron otros lenguajes ni entornos de ejecución. En segundo lugar, la detección se limitó a tres dimensiones estructurales —integridad de capas, dependencias cíclicas y concentración de dependencias salientes—, dejando por fuera otras manifestaciones de la deriva, como la erosión a nivel de servicios distribuidos o las desviaciones en los patrones de diseño. En tercer lugar, el grafo de dependencias se construyó a nivel de paquetes: las violaciones se clasifican con esa granularidad, si bien la evidencia que las acompaña desciende hasta la referencia concreta entre clases que las origina. Por último, el experimento controlado se realizó sobre un caso de estudio sintético, como aprobó el anteproyecto, y se complementó con validaciones sobre código real ---cuatro bibliotecas y tres aplicaciones en capas de código abierto, y treinta y cinco versiones publicadas de tres de ellas---, sin desplegar el mecanismo en un entorno productivo ni contar con la participación de sus equipos; la generalización a organizaciones con arquitecturas heterogéneas y la valoración de la utilidad del reporte por equipos reales quedaron, por tanto, planteadas como trabajo futuro y no como parte del alcance cumplido.

\section{Justificación del trabajo de grado}
Resolver el problema planteado aporta valor en varios frentes que se refuerzan entre sí. Desde la perspectiva económica, un mecanismo que detecte la deriva de forma temprana permite contener el crecimiento del costo de mantenimiento, cuya magnitud ilustra el cuarto del tiempo de desarrollo que, según \citet{Besker2019}, se pierde por la deuda técnica, y que la erosión arquitectónica tiende a agravar \citep{Pigazzini2021, Li2022}. Anticipar la detección equivale a corregir las desviaciones cuando todavía son baratas de revertir, en lugar de pagar su factura cuando ya comprometen la viabilidad del producto.

Desde el punto de vista técnico, la propuesta traslada a las PyMES una capacidad que hasta ahora se concentraba en organizaciones de mayor tamaño: la verificación continua de la conformidad arquitectónica. La evidencia disponible subraya la importancia de esta práctica para sostener la calidad de los sistemas a lo largo del tiempo \citep{Klein2022}, y herramientas como \textit{ArchUnit} ya demostraron que es posible expresar reglas arquitectónicas como pruebas ejecutables en Java \citep{ArchUnit}. Este trabajo aprovecha esa base conceptual, pero la orienta hacia un diagnóstico agregado y clasificatorio, más sencillo de interpretar para un equipo sin especialistas dedicados a la arquitectura.

En el plano social y organizacional, la solución empodera a los equipos de desarrollo, a los arquitectos y a los líderes técnicos de las PyMES, quienes pasan de depender de la intuición o de documentos obsoletos a disponer de una señal objetiva y reproducible sobre el estado estructural de sus proyectos. La propuesta se alinea, además, con las iniciativas de gobernanza técnica liviana que vienen cobrando fuerza en el ecosistema de las pequeñas organizaciones, en las que los marcos de control deben ser proporcionales a los recursos disponibles \citep{LAIG2025}.

\section{Metodología de la investigación}
La investigación se desarrolló, tal como se definió en el anteproyecto, como una investigación aplicada de enfoque predominantemente cuantitativo bajo la estrategia de \textit{Design Science Research} (DSR), adecuada cuando el propósito no es solo comprender un fenómeno, sino también diseñar y evaluar un artefacto que lo aborde \citep{Hevner2004, Peffers2007}. Para que el trabajo fuera alcanzable en el tiempo previsto, la estrategia se ejecutó mediante ciclos iterativos cortos, cada uno con un entregable verificable, organizados en las cinco fases que fijó el anteproyecto:

\begin{enumerate}
    \item \textbf{Delimitación del caso de estudio y catálogo de reglas.} Se seleccionó el ecosistema Java, se construyó el sistema base y se seleccionaron, a partir de la literatura, las tres \textit{Fitness Functions} del catálogo con su matriz de reglas (secciones~\ref{sec:seleccion} y~\ref{sec:escenarios}).
    \item \textbf{Diseño técnico del prototipo.} Se definieron el proceso de detección y sus entradas y salidas, la forma de declarar la arquitectura prevista, la de calcular e interpretar la Tasa de Conformidad y la arquitectura interna del prototipo, antes de programar (secciones~\ref{sec:proceso}, \ref{sec:declaracion}, \ref{sec:operacionalizacion}, \ref{sec:interpretacion} y~\ref{sec:diseno}).
    \item \textbf{Construcción del componente de extracción.} Se definieron las reglas con que se obtiene del código el grafo de dependencias entre paquetes y se implementó el módulo que las aplica (secciones~\ref{sec:extraccion} y~\ref{sec:implementacion}).
    \item \textbf{Construcción del motor de validación.} Sobre el grafo se implementaron la evaluación de las reglas, el cálculo de la CR, la clasificación por umbrales y el reporte de violaciones por tipo, ubicación y severidad (secciones~\ref{sec:evaluacion} e~\ref{sec:implementacion}).
    \item \textbf{Evaluación y análisis de resultados.} Se ejecutó el experimento controlado sobre la línea base y las versiones con derivas inducidas, y se contrastaron los resultados con la matriz de validación y con los indicadores comprometidos (capítulo~\ref{cap:evaluacion}). Como complemento, y sin ampliar el alcance de OBJ4, se repitió la prueba en entorno controlado sobre siete sistemas reales, inyectándoles derivas conocidas (sección~\ref{sec:sistemas}); se siguió la evolución de tres de ellos a lo largo de sus versiones publicadas para observar deriva real (sección~\ref{sec:evolucion}); se calibró el umbral de concentración con un procedimiento estadístico (sección~\ref{sec:umbral}), y se contrastó la extracción con una herramienta independiente (sección~\ref{sec:oraculo}).
\end{enumerate}

El ciclo no fue lineal, como la propia estrategia DSR prevé: la quinta fase reveló una limitación de la métrica de conformidad que obligó a volver a la cuarta, refinar el criterio de clasificación y reejecutar el experimento, de modo que el artefacto que se presenta es el producto de una segunda iteración y no de un único paso de construcción.

Metodológicamente, el trabajo articuló procedimientos de distinta naturaleza. La construcción del grafo de dependencias se apoyó en técnicas de análisis estático, que extraen la estructura del sistema sin necesidad de ejecutarlo. La evaluación de las propiedades estructurales se formalizó mediante un catálogo de \textit{Fitness Functions}, cada una de las cuales define sus propios criterios de verificación e incumplimiento. Finalmente, la agregación de resultados se resolvió con una métrica de conformidad acompañada de umbrales de clasificación, lo que permitió traducir los hallazgos individuales en un veredicto global sobre el estado del sistema. La validación se realizó mediante un experimento controlado que contrastó una implementación de referencia contra cuatro versiones deliberadamente alteradas, tres de ellas con una única deriva inducida para poder medir la sensibilidad de cada regla por separado, siguiendo las convenciones de diseño y reporte habituales en la investigación empírica en ingeniería de software \citep{Wohlin2012, RunesonHost2009}.

Las fases guardan con los objetivos específicos la correspondencia que estableció el anteproyecto, y que otorga trazabilidad al trabajo. La primera fase responde a OBJ1, la formulación del catálogo de \textit{Fitness Functions}; la segunda y la tercera, a OBJ2, la construcción de la base técnica del análisis estático y del grafo; la cuarta, a OBJ3, el motor de validación con la CR y sus umbrales; y la quinta, a OBJ4, el experimento controlado y el análisis de sus resultados. Los resultados de esa última fase retroalimentaron la cuarta, cerrando el ciclo iterativo que caracteriza a la investigación en diseño.

Se emplearon los cuatro instrumentos previstos. La \textbf{revisión bibliográfica focalizada} consolidó los criterios de deriva y la selección de las reglas (capítulo~\ref{cap:marco}). La \textbf{matriz de reglas arquitectónicas} registró, para cada regla, su racionalidad, su evidencia esperada y su criterio de incumplimiento, y es el catálogo de la tabla~\ref{tab:catalogo}. La \textbf{bitácora de desarrollo y ejecución} quedó en el historial de versiones del repositorio, que conserva cada iteración del prototipo, y en la evidencia que genera cada corrida, de modo que las decisiones y los resultados pueden reconstruirse; las decisiones con consecuencias conceptuales se registran además como \textit{Architecture Decision Records} en la carpeta \texttt{decisiones/} del repositorio, con su contexto, las alternativas descartadas y su respaldo (anexo~\ref{anexo:decisiones}). Las \textbf{tablas comparativas de resultados} registran el número de violaciones, la CR, el tipo de deriva y la clasificación de cada escenario frente a la línea base (sección~\ref{sec:resultados}).

\subsection{Ejecución del cronograma}
El anteproyecto planificó cinco meses, de mayo a septiembre de 2026, con 192 horas repartidas en 32 de revisión y definición de reglas, 80 de diseño e implementación, 42 de evaluación y ajustes y 38 de documentación y cierre, y previó al cierre de cada mes una revisión con el director para contrastar horas planeadas y reales. Los entregables de cada mes se produjeron en el orden previsto: el protocolo y el catálogo inicial (capítulo~\ref{cap:desarrollo}, primera parte), el diseño del prototipo y la extracción (secciones~\ref{sec:extraccion} y~\ref{sec:diseno}), las reglas ejecutables, la detección de ciclos con el cálculo de la CR y el reporte, y el estudio de caso con su análisis comparativo (capítulo~\ref{cap:evaluacion}). El trabajo se extendió hasta octubre de 2026, mes en que se añadieron las validaciones complementarias sobre código real y se consolidó el documento.
% TODO ENTREGA: si el autor registró las horas reales, añadirlas aquí junto a las planeadas.

Dos desviaciones se declaran en lugar de omitirse. La primera es que \textbf{las revisiones mensuales con el director no se realizaron} durante la ejecución, de modo que la bitácora no contiene sus observaciones y la validación con él se concentra en la revisión de este documento. La segunda es que, a falta de esas revisiones, el contraste mensual entre horas planeadas y reales que preveía la metodología no quedó documentado; el historial de versiones permite, en cambio, fechar cada entregable. Ninguna de las dos afecta a los resultados técnicos, pero ambas restan la supervisión externa que el anteproyecto había previsto como control de calidad del proceso, y por eso se registran como decisión en el anexo~\ref{anexo:decisiones}.
% TODO ENTREGA: ajustar este párrafo si hubo revisiones con el director antes de radicar.

\section{Resultados obtenidos}
El mecanismo desarrollado logró construir, a partir del código compilado, un grafo de dependencias entre paquetes y ejecutar sobre él las tres \textit{Fitness Functions} previstas. Sobre la implementación de referencia obtuvo una Tasa de Conformidad del cien por ciento, sin un solo hallazgo, y la clasificó como \textit{aceptable}. Sobre las cuatro versiones alteradas detectó la totalidad de las seis derivas inducidas —rupturas de la integridad de capas, dependencias cíclicas y concentración excesiva de dependencias salientes—, superando el criterio de éxito del ochenta por ciento fijado en el anteproyecto. En el escenario que combina las tres derivas, la conformidad descendió al cincuenta y siete por ciento y el sistema fue clasificado como \textit{no conforme}.

La evaluación aportó, además, un hallazgo que no se había anticipado en su magnitud. Los escenarios construidos con una única deriva revelaron que la Tasa de Conformidad, al ser una proporción no ponderada, diluye las violaciones aisladas: en dos de esos escenarios el mecanismo detectó correctamente la desviación pero la métrica se mantuvo por encima del umbral de aceptación, de modo que el sistema quedaba clasificado como aceptable pese a contener una violación real. Ese resultado motivó un refinamiento del criterio de clasificación —que incorpora la gravedad del peor hallazgo junto a la extensión que mide la métrica, sin alterar su cálculo— con el cual los cinco escenarios pasaron a clasificarse conforme a lo esperado.

Al margen del experimento controlado, el mecanismo se aplicó sobre una biblioteca de terceros de casi mil clases para observar su comportamiento ante código que nadie preparó para él. El análisis completo tardó poco más de un segundo y la regla de integridad de capas no produjo un solo falso positivo. Reveló, en cambio, una segunda limitación: sobre código real, una regla puramente topológica confunde un componente deteriorado con uno que concentra dependencias por el papel que desempeña, hasta el punto de que dos tercios de sus hallazgos resultaron ser ruido. La respuesta fue incorporar exenciones declaradas y registradas en el reporte, que eliminan ese ruido sin silenciar los hallazgos genuinos.

Finalmente, para comprobar que lo observado en el caso de estudio se sostiene fuera de él, se inyectaron derivas conocidas en siete sistemas Java reales ---cuatro bibliotecas de entre cuatrocientas y dos mil clases y tres aplicaciones en capas cuya arquitectura documentan sus propios proyectos---. El mecanismo detectó las treinta y una derivas inyectadas sin producir un solo hallazgo espurio y en menos de un segundo por sistema, y la dilución de la Tasa de Conformidad reapareció con mayor claridad: en el sistema más conforme, una deriva aislada la dejaba por encima del noventa y siete por ciento. Sobre treinta y cinco versiones publicadas de tres bibliotecas, el mecanismo encontró además deriva real, que nadie inyectó: una dependencia nueva de la introspección hacia la deserialización en Jackson Databind 2.19.0 y una concentración nueva en Commons Lang 3.21.0, mientras que Guava conservó su arquitectura a lo largo de doce versiones publicadas. Una herramienta independiente, JDepend, extrajo el mismo grafo salvo veintiuna aristas, todas explicadas por las reglas de extracción documentadas, y las mismas reglas escritas con \textit{ArchUnit} detectaron lo mismo con más del doble de especificación y sin veredicto agregado. Los capítulos de desarrollo y evaluación detallan los resultados y el razonamiento que condujo a cada ajuste.

\section{Estructura del documento}
El presente documento se organiza en seis capítulos y tres anexos. El capítulo~\ref{cap:marco}, \textit{Marco de referencia}, expone las bases teóricas sobre deriva arquitectónica, \textit{Fitness Functions}, análisis estático, grafos de dependencia y métricas de acoplamiento, junto con el estado del arte. El capítulo~\ref{cap:desarrollo}, \textit{Desarrollo del proyecto}, es el núcleo del trabajo y se divide en dos partes: la primera explica el proceso de detección ---qué entidades intervienen, de dónde proceden y por qué se seleccionaron las \textit{Fitness Functions}, cómo se declara la arquitectura prevista, con qué reglas se extrae el grafo, cómo se obtiene la deriva, en qué unidades se mide y cómo se interpreta---, y la segunda describe el prototipo que lo instrumenta y los escenarios del caso de estudio. El capítulo~\ref{cap:evaluacion}, \textit{Evaluación}, presenta el experimento controlado, su interpretación, el refinamiento que motivó, la aplicación a sistemas reales y las amenazas a la validez. El capítulo \textit{Conclusiones} responde a los objetivos, precisa la aplicación del mecanismo en la industria y traza las líneas de trabajo futuro, y el capítulo~\ref{cap:riesgos}, \textit{Riesgos éticos}, examina las salvaguardas asociadas a su uso. El anexo~\ref{anexo:repro} reúne las instrucciones para reproducir todos los resultados, el anexo~\ref{anexo:asistente} documenta una exploración del asistente de catálogo con modelos de lenguaje reales, y el anexo~\ref{anexo:decisiones} resume el registro de decisiones del trabajo.